\documentclass[10pt,a4paper]{amsart}
\usepackage[margin=19mm]{geometry}
\usepackage{amsmath,amssymb,mathtools}
\usepackage[scr=rsfs,cal=euler]{mathalfa}
\usepackage{xfrac}
\usepackage[unicode,hidelinks,bookmarks=false]{hyperref}
\newcommand{\ie}{i.e.\ }
\newcommand{\cf}{cf.\ }
\newcommand{\resp}{respectively\ }
\newcommand{\Subsection}{\subsection}
\providecommand{\nobreakdash}{\nobreak\hbox{-}}
\setlength{\emergencystretch}{2em}
\multlinegap0pt
\usepackage{tikz}
\usepackage{amssymb}
\usepackage[all]{xy}
\usepackage{stackengine}
\usepackage{enumerate}
\usepackage{xfrac}
\newtheorem{proposition}{Proposition}
\newcommand{\K}{ {\Bbbk} }
 \newcommand{\KP}{{\Bbbk[X]/(P_1)}}
\newcommand{\KPP}{{\Bbbk[X]/(P_2)}}
 \newcommand{\KPPn}{{\Bbbk[X]/(P_2^n)}}
 \newcommand{\KPn}{{\Bbbk[X]/(P_1^n)}}
\title{on some local rings}
\author{Mohamad Maassarani}
\date{Source: arXiv:2512.19197v1 (CC BY 4.0)}
\begin{document}
\maketitle

\noindent\emph{Source and licence.} on some local rings. Version arXiv:2512.19197v1 (CC BY 4.0), distributed by arXiv under the Creative Commons Attribution 4.0 International licence, http://creativecommons.org/licenses/by/4.0/, as recorded in the arXiv licence metadata for this exact version and shown on the abstract page https://arxiv.org/abs/2512.19197v1. Full licence notice: \texttt{licenses/arxiv-2512.19197-CC-BY-4.0.txt}.\par\smallskip

\noindent\textit{Editorial scope.} Counted unit: the article's proposition proposition (source0.tex lines 217--223). The article's complete original proof of it is delivered with it (source0.tex lines 224--231, one \texttt{proof} environment of 8 lines). No mathematics is written, repaired or re-derived: the statement and the proof are the article's own lines, in the article's own notation, and no retained line is substituted or re-flowed.  Retained uncounted as the source-local context the proof needs: lines 192--201.  Nothing was omitted from any retained span, so the package carries no omission record.  No retained line contains a citation, so the wrapper carries no bibliography and no bibliography entry.  No retained line contains a figure environment, an image-inclusion command or drawing code, so no image is delivered and the package declares no asset dependency.  Editorial formatting only, in addition: an AMS \texttt{amsart} preamble that restates the article's own declarations and macros used by the retained lines, namely the environments \texttt{proposition} on separate counters, together with the standard packages those lines need.  The retained environments print in this document as Proposition 1 in order of appearance, whereas the article prints the same items with the numbering of its own full document.  The complete native source of the article as distributed in the arXiv e-print for this exact version is bundled unchanged as \texttt{sources/191451/source0.tex}.
\begin{proposition}\label{P1}
If $f: \KP\to \KPP$ is a ring isomorphism stabilizing $\K$ then :
\begin{itemize}
\item[1)] The degree of $P_1$ is equal to the degree of $P_2$.
\item[2)] There exist a unique polynomial $Q_f\in \K[X]$ of degree less than the degree of $P_1$ (the degree of $P_2$) and greater or equal to $1$ such that $f$ is induced by the ring morphism stabilizing $\K$ $f_X: \K[X]\to \K[X]$ defined by $X\mapsto Q_f  $ and $\sigma_{f_X}=\sigma_f$ i.e. $P\mapsto \sigma_{f}^X(P)\circ Q_f$, where $\sigma_f^X$ is as in the previous proposition.
\item[3)] $\sigma_f^X(P_1)\circ Q_f=S_fP_2$ for a given $S_f \in \K[X]$.
\item[4)] For $P\in K[X]$, if $\sigma_f^X(P)\circ Q_f= SP_2$ for some $S\in \K[X]$ then $P=RP_1$ for some $R\in \K[X]$.  
\item[5)] The morphism $f_X$ maps $(P_1^n)$ into $(P_2^n)$ and hence induces a ring morphism stabilizing $\K$ :  $f_{X,n}:\KPn\to \KPPn$ induced by $P\mapsto \sigma_f^X(P)\circ Q_f$.
\end{itemize}
\end{proposition}

\begin{proposition}
Let $f: \KP\to \KPP$ be a ring isomorphism stabilizing $\K$.
\begin{itemize}
\item[1)] If $\alpha$ is a root of $P_2$ then $Q_f(\alpha)$ is a root of $\sigma_f^X(P_1)$.
\item[2)] We have a bijection $\{\text{roots of } P_2\}\rightarrow \{\text{roots of } \sigma_f^X(P_1)\}$ given by $\alpha \mapsto Q_f(\alpha)$.
\end{itemize}
\end{proposition}

\begin{proof}
$1)$ follows from the equation $\sigma_f^X(P_1)\circ Q= S_f P_2$ of proposition \ref{P1} and it follows from $1)$ that we have a map $g: \{\text{roots of } P_2\}\rightarrow \{\text{roots of } \sigma_f^X(P_1)\}$ given by $\alpha \mapsto Q_f(\alpha)$. To prove $2)$ we will define an inverse to $g$. Applying $3)$ of propsition \ref{P1} to $f^{-1}$ we get that :
$$(\sigma_f^X)^{-1}(P_2)\circ Q_{f^{-1}}=S_{f^-1}P_1,$$
and hence $$P_2 \circ \sigma_f^X(Q_{f^{-1}})=\sigma_f^X (S_{f^{-1}})\sigma_f^X (P_1).$$
Therefore we have a well defined map  $h:\{\text{roots of } \sigma_f^X(P_1)\} \to \{\text{roots of } P_2\}$ given by $\alpha\mapsto \sigma_f^X(Q_{f^{-1}})(\alpha)$. We will prove that $h$ and $g$ are inverse to each other.  Since $f$ and $f^{-1}$ are inverse to each other, we have :
$$ \sigma_f^X(Q_{f^{-1}})\circ Q_f = X +S_2P_2 \quad \text{and} \quad (\sigma_f^X)^{-1}(Q_f)\circ Q_{f^{-1}}=X+S_1 P_1,$$ 
for some $S_1,S_2 \in \K[X]$. The first equation shows that $hg$ is the identity of $\{\text{roots of } P_2\}$. Composing the second equation by $\sigma_f^X$ we obtain that $gh$ is the identity of $\{\text{roots of } \sigma_f^X(P_1)\}$. We have proved $2)$.
\end{proof}

\end{document}
